% !TeX root = closed-systems-from-comparison-completeness.tex
% ============================================================
\chapter{Nonperturbative Adiabatic Invariants from Closed-System Transport}
\label{app:fusion}
% ============================================================
% ============================================================
\begin{chapterorientation}
\orientationitem{Input}{Closed-system transport, the stabilized quadratic scalar channel, and the conditional realization distinctions of the main text.}
\orientationitem{Role}{This appendix records a downstream confinement application of the transport package.}
\orientationitem{Output}{A nonperturbative invariant package, with intrinsic scalar content separated from perturbative and magnetic-confinement readings.}
\end{chapterorientation}
\section{Purpose and logical status}
\label{sec:fusion-purpose}
Under the standing principle of closed-world admissibility
(\cref{sp:closed-world}), this appendix records a downstream
confinement application of \cref{sp:compactness,sp:transport} through
the stabilized quadratic scalar channel.
As the terminal step of the manuscript, it isolates the theorem-level
intrinsic scalar carried by any confinement system realized as a closed
comparison world and then separates that scalar from its later
perturbative and magnetic-confinement readings.
% ============================================================
The purpose of this appendix is to present a downstream application
template for the closed relational stack developed in the main text.
The chosen downstream model is magnetic confinement.
The conclusion at theorem level is that, once the transport stack and
the inherited realized scalar-channel package of
\cref{thm:interface-stabilized,thm:born,thm:quadratic-scalar} are in
force, one obtains a scalar
\[
	\mathfrak s_{\mathrm{fus}}:U\to\mathbb R
\]
which is unique up to nonzero normalization and invariant under the
intrinsic evolution of the confinement system.
This appendix is entirely downstream.
It introduces no new foundational principle, assumes no new dynamical
axiom locally, and does not reprove the closed-system stack.
Its role is to show that once the classification, quotient,
transport, quadratic-carrier, and realized scalar-channel theorems are
available, the intrinsic scalar
\[
	\mathfrak s_{\mathrm{fus}}=Q\circ \kappa
\]
follows formally.
The operative chain is:
\begin{align*}
	(U,\mathcal C)
	 & \Longrightarrow U\cong X_A\times X_B                    \\
	 & \Longrightarrow \pi:X\to \Phys=X/G                      \\
	 & \Longrightarrow F^1\supset F^2\supset F^3\supset \cdots \\
	 & \Longrightarrow \mathcal K_\infty\subset F^2/F^3        \\
	 & \Longrightarrow Q:\mathcal K_\infty\to\mathbb R       \\
	 & \Longrightarrow \mathfrak s_{\mathrm{fus}}=Q\circ \kappa .
\end{align*}
The point is exact at the level of the intrinsic scalar.
The quadratic scalar channel isolated in the main text is therefore not
left as a formal carrier datum.
In the confinement setting it furnishes a canonical nonperturbative
transport scalar; later perturbative matching and physical
interpretation are treated separately below.
\begin{theorem}[Main appendix theorem]
	\label{thm:fusion-main-statement}
	Under the standing principle \cref{sp:closed-world}, let
	\((U,\mathcal C)\) be the closed comparison world associated to a
	magnetic confinement system.
	Assume the structural inputs listed in
	\cref{sec:fusion-structural-inputs}.
	Then there exists a scalar-valued function
	\[
		\mathfrak s_{\mathrm{fus}}:U\to\mathbb R
	\]
	such that:
	\begin{enumerate}[label=\textup{(\alph*)}]
		\item
		      \(\mathfrak s_{\mathrm{fus}}\) is preserved by every intrinsic evolution
		      \[
			      \Phi_t\in \Aut(U,\mathcal C),
		      \]
		      that is,
		      \[
			      \mathfrak s_{\mathrm{fus}}(\Phi_t(u))=\mathfrak s_{\mathrm{fus}}(u)
			      \qquad
			      \text{for all }u\in U,\ t;
		      \]
		\item
		      \(\mathfrak s_{\mathrm{fus}}\) is unique up to multiplication by a nonzero real scalar.
	\end{enumerate}
\end{theorem}
\begin{proof}
	Existence is proved in \cref{thm:fusion-existence}, invariance under
	intrinsic evolution in \cref{thm:fusion-invariance}, and uniqueness up
	to nonzero scaling in \cref{thm:fusion-unique}.
	Combining these three statements yields items
	\textup{(a)}--\textup{(b)}.
\end{proof}
% ============================================================
\section{Structural inputs from the main text}
\label{sec:fusion-structural-inputs}
% ============================================================
The appendix uses the transport/carrier results already
established in the main text together with the realized scalar-channel
package already isolated there.
No new primitive axiom is introduced locally.
The required inputs are:
\begin{enumerate}[label=\textup{(\roman*)}]
	\item
	      the closed-world quotient backbone and structural completion of
	      admissible state content
	      \(\bigl(\cref{thm:part1-structural-completion}\bigr)\);
	\item
	      the lift and transport-obstruction architecture, including the
	      canonical morphism-locus carrier and its first visible quadratic
	      layer
	      \(\bigl(\cref{thm:lift-classification-main,thm:triangle-degree2}\bigr)\);
	\item
	      stabilization of the intrinsic obstruction channel and its
	      realization interface
	      \(\bigl(\cref{thm:interface-stabilized}\bigr)\);
	\item
	      existence and one-dimensional rigidity of the scalar channel on the
	      realized degree--\(2\) channel attached to the stabilized
	      quadratic carrier, together with its probabilistic Born
	      normalization under the inherited second-jet interface
	      hypothesis
	      \(\bigl(\cref{thm:quadratic-scalar,thm:born}\bigr)\).
\end{enumerate}
These inputs are exactly the data propagated in the remainder of the
appendix to produce the intrinsic confinement scalar
\(\mathfrak s_{\mathrm{fus}}=Q\circ\kappa\).
Any perturbative identification or magnetic-confinement reading is kept
as a downstream interpretation layer on top of that scalar.
% ============================================================
\section{Transport filtration and the universal quadratic quotient}
\label{sec:fusion-universal}
% ============================================================
Let
\[
	N:=\Ker\!\bigl(\Aut(U,\mathcal C)\to \Aut(X_A)\times\Aut(X_B)\bigr),
	\qquad
	F^1:=N,
	\qquad
	F^{m+1}:=[F^1,F^m].
\]
\begin{definition}[Visible degree--\(1\) transport module]
	\[
		V:=F^1/F^2.
	\]
\end{definition}
% ------------------------------------------------------------
\subsection*{Commutator structure of the transport filtration}
% ------------------------------------------------------------
\begin{lemma}
	\label{lem:fusion-commutator}
	For all integers \(i,j\ge 1\),
	\[
		[F^i,F^j]\subseteq F^{i+j}.
	\]
	In particular, each quotient \(F^m/F^{m+1}\) is abelian.
\end{lemma}
\begin{proof}
	We proceed in a sequence of logically independent steps.
	\medskip
	\noindent
	\textbf{Step 1: Reduction to a universal group-theoretic statement.}
	The filtration is defined inductively by
	\[
		F^{m+1}=[F^1,F^m].
	\]
	Thus \(F^\bullet\) is the lower central series of the group \(F^1\).
	It suffices to prove that for any group \(G\) with lower central series
	\[
		\gamma_1(G)=G,\qquad \gamma_{m+1}(G)=[G,\gamma_m(G)],
	\]
	one has
	\[
		[\gamma_i(G),\gamma_j(G)]\subseteq \gamma_{i+j}(G).
	\]
	We therefore prove the statement in this generality.
	\medskip
	\noindent
	\textbf{Step 2: Induction scheme.}
	Define
	\[
		P(n):\quad
		[\gamma_i(G),\gamma_j(G)]\subseteq \gamma_{i+j}(G)
		\quad\text{for all }i+j\le n.
	\]
	We prove \(P(n)\) for all \(n\ge2\) by induction.
	\medskip
	\noindent
	\textbf{Step 3: Base case.}
	For \(n=2\), the only possibility is \(i=j=1\).
	Then
	\[
		[\gamma_1(G),\gamma_1(G)]
		=
		[G,G]
		=
		\gamma_2(G),
	\]
	so the statement holds.
	\medskip
	\noindent
	\textbf{Step 4: Inductive step.}
	Assume \(P(n-1)\).
	Let \(i+j=n\).
	Without loss of generality, assume \(i\ge2\).
	Then
	\[
		\gamma_i(G)=[G,\gamma_{i-1}(G)].
	\]
	Thus
	\[
		[\gamma_i(G),\gamma_j(G)]
		=
		[[G,\gamma_{i-1}(G)],\gamma_j(G)].
	\]
	We now invoke the Hall--Witt identity (three-subgroup identity):
	for subgroups \(A,B,C\),
	\[
		[[A,B],C]\subseteq [A,[B,C]]\cdot [B,[A,C]].
	\]
	We use the standard group-theoretic form of this identity
	\cite{Hall1959}.
	Apply with
	\[
		A=G,\quad B=\gamma_{i-1}(G),\quad C=\gamma_j(G).
	\]
	Then
	\[
		[\gamma_i(G),\gamma_j(G)]
		\subseteq
		[G,[\gamma_{i-1},\gamma_j]]
		\cdot
		[\gamma_{i-1},[G,\gamma_j]].
	\]
	By the induction hypothesis,
	\[
		[\gamma_{i-1},\gamma_j]\subseteq \gamma_{i+j-1}(G),
	\]
	and by definition,
	\[
		[G,\gamma_j]=\gamma_{j+1}(G).
	\]
	Therefore,
	\[
		[G,[\gamma_{i-1},\gamma_j]]
		\subseteq
		[G,\gamma_{i+j-1}]
		=
		\gamma_{i+j}(G),
	\]
	and
	\[
		[\gamma_{i-1},[G,\gamma_j]]
		=
		[\gamma_{i-1},\gamma_{j+1}]
		\subseteq
		\gamma_{(i-1)+(j+1)}(G)
		=
		\gamma_{i+j}(G).
	\]
	Thus
	\[
		[\gamma_i(G),\gamma_j(G)]
		\subseteq
		\gamma_{i+j}(G).
	\]
	\medskip
	\noindent
	\textbf{Step 5: Abelianity of graded quotients.}
	Taking \(i=j=m\),
	\[
		[\gamma_m(G),\gamma_m(G)]\subseteq \gamma_{2m}(G).
	\]
	Since \(2m\ge m+1\),
	\[
		\gamma_{2m}(G)\subseteq \gamma_{m+1}(G),
	\]
	hence
	\[
		[\gamma_m(G),\gamma_m(G)]\subseteq \gamma_{m+1}(G).
	\]
	Therefore the quotient
	\[
		\gamma_m(G)/\gamma_{m+1}(G)
	\]
	is abelian.
	\medskip
	\noindent
	\textbf{Step 6: Specialization.}
	Applying this to \(G=F^1\), we obtain
	\[
		[F^i,F^j]\subseteq F^{i+j},
	\]
	as required.
\end{proof}
% ------------------------------------------------------------
\subsection*{Universal quadratic quotient}
% ------------------------------------------------------------
\begin{theorem}[Universal quadratic quotient]
	There exists a canonical surjective homomorphism
	\[
		q_{\wedge}:\Lambda^2(V)\twoheadrightarrow F^2/F^3
	\]
	characterized by
	\[
		q_{\wedge}(\bar x\wedge \bar y)=[x,y]\bmod F^3.
	\]
\end{theorem}
\begin{proof}
	We divide the proof into five logically independent parts.
	\medskip
	\noindent
	\textbf{Step 1: Definition of the pairing.}
	Let \(\bar x,\bar y\in V=F^1/F^2\).
	Choose representatives \(x,y\in F^1\).
	Define
	\[
		\beta(\bar x,\bar y):=[x,y]\bmod F^3.
	\]
	We must show that \(\beta\) is:
	\begin{enumerate}
		\item well defined;
		\item bilinear;
		\item alternating;
		\item universal;
		\item surjective.
	\end{enumerate}
	\medskip
	\noindent
	\textbf{Step 2: Independence of representatives.}
	Suppose
	\[
		x'=xf,\quad y'=yg,
		\quad f,g\in F^2.
	\]
	We compute
	\[
		[x',y']=[xf,yg].
	\]
	Using the standard commutator expansion
	(valid in any group; see \cite{Hall1959}),
	\[
		[ab,cd]
		=
		[a,d]\cdot[a,c]^d\cdot[b,d]^c\cdot[b,c],
	\]
	and simplifying modulo \(F^3\), we obtain the classical reduction:
	\[
		[xf,yg]
		=
		[x,y]\cdot[x,g]\cdot[f,y]\cdot[f,g]
		\quad \mod F^3.
	\]
	Now:
	- \([x,g]\in [F^1,F^2]=F^3\),
	- \([f,y]\in [F^2,F^1]=F^3\),
	- \([f,g]\in [F^2,F^2]\subseteq F^4\subseteq F^3\).
	Hence
	\[
		[xf,yg]\equiv [x,y]\pmod{F^3}.
	\]
	Thus \(\beta\) is well defined.
	\medskip
	\noindent
	\textbf{Step 3: Bilinearity.}
	Let \(\bar x,\bar x',\bar y\in V\), with lifts \(x,x',y\in F^1\).
	We compute:
	\[
		[xx',y]
		=
		[x,y]\cdot[x,y,x']\cdot[x',y],
	\]
	where
	\[
		[x,y,x']:=([[x,y],x']).
	\]
	Since \([x,y]\in F^2\),
	\[
		[x,y,x']\in [F^2,F^1]=F^3.
	\]
	Thus modulo \(F^3\),
	\[
		[xx',y]\equiv [x,y]\cdot[x',y].
	\]
	Passing to the quotient \(F^2/F^3\), multiplication becomes addition, hence
	\[
		\beta(\bar x+\bar x',\bar y)
		=
		\beta(\bar x,\bar y)+\beta(\bar x',\bar y).
	\]
	The same argument applies in the second variable.
	\medskip
	\noindent
	\textbf{Step 4: Alternation.}
	We have
	\[
		[x,x]=e \Rightarrow \beta(\bar x,\bar x)=0.
	\]
	Also
	\[
		[y,x]=[x,y]^{-1},
	\]
	which corresponds to negation in the abelian group \(F^2/F^3\).
	Thus \(\beta\) is alternating.
	\medskip
	\noindent
	\textbf{Step 5: Factorization.}
	By the universal property of the exterior square,
	every alternating bilinear map
	\[
		V\times V \to F^2/F^3
	\]
	factors uniquely through
	\[
		\Lambda^2(V).
	\]
	Thus there exists a unique homomorphism
	\[
		q_{\wedge}:\Lambda^2(V)\to F^2/F^3
	\]
	with
	\[
		q_{\wedge}(\bar x\wedge \bar y)=\beta(\bar x,\bar y).
	\]
	\medskip
	\noindent
	\textbf{Step 6: Surjectivity.}
	By definition,
	\[
		F^2=[F^1,F^1].
	\]
	Thus every element of \(F^2\) is a finite product of commutators:
	\[
		z=\prod_i [x_i,y_i].
	\]
	Passing to \(F^2/F^3\), the group becomes abelian, so
	\[
		z\equiv \sum_i [x_i,y_i].
	\]
	Hence every element lies in the image of \(q_{\wedge}\).
	\medskip
	\noindent
	\textbf{Conclusion.}
	Therefore \(q_{\wedge}\) is a well-defined surjective homomorphism.
\end{proof}
% ------------------------------------------------------------
\begin{definition}[Quadratic relation subgroup]
	Define
	\[
		R_2:=\Ker(q_{\wedge})\subseteq \Lambda^2(V).
	\]
	Then there is an exact sequence
	\[
		0\longrightarrow R_2
		\longrightarrow \Lambda^2(V)
		\xrightarrow{q_{\wedge}}
		F^2/F^3
		\longrightarrow 0.
	\]
\end{definition}
\begin{remark}
	This identifies the universal degree--\(2\) commutator object.
	The invariant theorem does not require solving the full relation problem
	on \(\Lambda^2(V)\).
	By \cref{thm:kernel-annihilation-upstream}, every admissible degree--\(2\)
	scalarization already vanishes on \(R_2\), and therefore factors uniquely
	through the quadratic carrier. Thus the scalar theory of the appendix is
	intrinsic to the quadratic carrier \(F^2/F^3\) and not to any choice of universal lift.
\end{remark}
% ============================================================
\section{The stabilized quadratic carrier and its realized scalar channel}
\label{sec:fusion-K}
% ============================================================
\Cref{sec:fusion-universal} constructs the universal degree--\(2\)
quotient.
We now pass to the stabilized carrier that actually controls the
invariant.
\begin{proposition}
	The compatible degree--\(2\) transport defects define the stabilized
	quadratic carrier
	\[
		K=\mathcal K_\infty
	\]
	of \cref{thm:interface-stabilized}.
\end{proposition}
\begin{proof}
	We make explicit the logical content of the cited theorem in the notation
	of the present appendix.
	By construction, the transport package of the main text assigns to each
	finite stage of the refinement tower a degree--\(2\) carrier
	\[
		\mathcal K_k
		\subseteq
		F^2K_{p_k}^{(k)}/F^3K_{p_k}^{(k)}
	\]
	generated by the quadratic square classes arising from the intrinsic
	transport defect at that stage.
	The refinement morphisms induce compatible transition maps
	\[
		\mathcal K_{k+1}\longrightarrow \mathcal K_k
	\]
	because the defect construction commutes with refinement; this is the
	content of the compatibility statements proved upstream and assembled in
	\cref{thm:interface-stabilized}.
	A compatible family
	\[
		(\xi_k)_k,\qquad \xi_k\in \mathcal K_k,
	\]
	is by definition an element of the inverse limit
	\[
		\varprojlim_k \mathcal K_k.
	\]
	Conversely, every element of the inverse limit is, by definition, such a
	compatible family. The cited theorem identifies this inverse-limit object
	as the stabilized degree--\(2\) transport carrier and denotes it by
	\[
		\mathcal K_\infty.
	\]
	Therefore the stabilized carrier relevant for the present appendix is
	precisely
	\[
		K:=\mathcal K_\infty.
	\]
	No additional structure is introduced here; the proposition is exactly
	the translation of \cref{thm:interface-stabilized} into the present
	notation.
\end{proof}
\begin{proposition}
	\label{prop:fusion-curvature}
	Under the inherited second-jet faithfulness condition and compatible
	smooth realization, nonzero stabilized square classes in the
	stabilized quadratic carrier are equivalent to nonzero realized
	curvature on the realized degree--\(2\) channel.
\end{proposition}
\begin{proof}
	The statement has two parts: first, the passage from intrinsic quadratic
	transport defect to the stabilized carrier; second, the realized
	interface criterion carried by compatible smooth realization.
	By \cref{thm:interface-curvature}, if
	\[
		\mathcal A_{ij}(p_\infty)\in \mathcal K_\infty
	\]
	denotes the stabilized quadratic square class determined by the
	\((i,j)\)-transport square at the inverse-limit basepoint \(p_\infty\),
	then under the second-jet comparison map
	\[
		\Jtwo:\mathcal K_\infty\to \End(V),
		\qquad
		V:=T_{\iota(p_\infty)}M,
	\]
	nonzero stabilized square classes are equivalent to nonzero realized
	curvature on the realized degree--\(2\) channel.
	By \cref{thm:connection-smooth-realization}, the same compatible smooth
	realization supplies the realized affine-connection and metric branch on
	which this curvature is read.  Thus the appendix inherits the
	main-text intrinsic--smooth interface exactly: the stabilized quadratic carrier feeds the
	confinement application through the inherited second-jet criterion
	relating its nonzero stabilized square classes to nonzero realized
	curvature, not through a direct whole-carrier curvature identity.
\end{proof}
\begin{proposition}[Realized scalar channel on the stabilized carrier]
	\label{prop:fusion-scalar}
	Under smooth realization and the inherited second-jet interface
	hypothesis, the stabilized quadratic carrier admits a nonzero scalar
	channel
	\[
		Q:K\to\mathbb R
	\]
	read from the realized degree--\(2\) channel attached to that carrier,
	unique up to multiplication by a nonzero real scalar.
\end{proposition}
\begin{proof}
	By \cref{prop:fusion-curvature}, under smooth realization and the
	inherited second-jet faithfulness condition, the stabilized quadratic
	carrier
	\[
		K=\mathcal K_\infty
	\]
	is read through the realized curvature carrier. The scalarization
	problem used in this appendix is therefore exactly the scalarization
	problem on the realized degree--\(2\) channel addressed in the main text.
	Now \cref{lem:K-scalar-unique,lem:branch-weight-1d} state that the space
	of admissible scalarizations of that realized degree--\(2\) channel is
	one-dimensional. Concretely, this means that if
	\[
		W_1,W_2:K\to\mathbb R
	\]
	are admissible scalarizations, then there exists
	\[
		c\in\mathbb R
	\]
	such that
	\[
		W_1=c\,W_2.
	\]
	In particular, the scalarization space is nonzero, so there exists at
	least one nontrivial scalarization. Fix any such nonzero scalarization
	and denote it by
	\[
		Q:K\to\mathbb R.
	\]
	If \(W:K\to\mathbb R\) is any other admissible scalarization, then the
	one-dimensionality assertion implies that there exists
	\[
		c\in\mathbb R
	\]
	with
	\[
		W=cQ.
	\]
	If \(W\neq 0\), then necessarily \(c\neq 0\). Hence every nonzero
	admissible scalarization is a nonzero scalar multiple of \(Q\).
	This is the scalar channel used in the remainder of the appendix.
\end{proof}
% ------------------------------------------------------------
% ------------------------------------------------------------
\subsection*{Canonical quadratic defect section}
% ------------------------------------------------------------
The construction of the confinement invariant uses a canonical map
\[
	\kappa:U\to K
\]
from the state space to the stabilized quadratic carrier.
To make the downstream scalarization completely rigid, we isolate the
construction in three stages:
existence, basepoint-independence, and functorial uniqueness.
\begin{theorem}[Canonical quadratic defect section]
	\label{thm:fusion-kappa}
	Under the standing principle \cref{sp:closed-world}, let
	\((U,\mathcal C)\) be a closed comparison world, let
	\[
		G:=\Aut(U,\mathcal C),
	\]
	and let
	\[
		K=\mathcal K_\infty
	\]
	be the stabilized quadratic transport carrier obtained from the
	refinement-compatible degree--\(2\) transport system of
	\cref{ch:interface}.
	In the local witness-selection setup, assume that for each state
	\(u\in U\), the refinement tower admits a
	compatible family of local transport witnesses centered at \(u\), so
	that at each finite stage \(k\) one has a local witness network
	\[
		\mathcal N_k(u)=(V_k(u),E_k(u)),
	\]
	a chosen admissible basepoint
	\[
		p_k(u)\in V_k(u),
	\]
	and associated quadratic square classes
	\[
		\mathcal A_{ij}^{(k)}(p_k(u))
		\in
		F^2K_{p_k(u)}^{(k)}/F^3K_{p_k(u)}^{(k)}.
	\]
	Assume moreover that these classes are compatible under the refinement
	bonding maps.
	Then the family
	\[
		\bigl(\mathcal A_{ij}^{(k)}(p_k(u))\bigr)_k
	\]
	defines a canonical element
	\[
		\kappa(u)\in K.
	\]
	Hence there exists a canonical map
	\[
		\kappa:U\to K,
		\qquad
		u\longmapsto \kappa(u),
	\]
	called the quadratic defect section.
	Furthermore, for every
	\[
		g\in G
		\qquad\text{and}\qquad
		u\in U,
	\]
	one has
	\[
		\kappa(g\cdot u)=g_*\kappa(u),
	\]
	where
	\[
		g_*:K\to K
	\]
	is the induced automorphism of the stabilized quadratic carrier.
\end{theorem}
\begin{proof}
	We divide the argument into six steps.
	\medskip
	\noindent
	\textbf{Step 1: Finite-stage quadratic defect classes attached to a state.}
	Fix
	\[
		u\in U.
	\]
	By hypothesis, at each refinement stage \(k\), the state \(u\) determines
	a local transport witness system with basepoint
	\[
		p_k(u)\in V_k(u),
	\]
	and therefore a based defect homomorphism
	\[
		\delta_{p_k(u)}^{(k)}:
		\Omega_{p_k(u)}(\mathcal N_k(u))\to K_{p_k(u)}^{(k)}.
	\]
	For each ordered pair of local transport directions \((i,j)\), let
	\[
		\square_{ij}:=e_ie_je_i^{-1}e_j^{-1}
	\]
	be the corresponding commutator square loop based at \(p_k(u)\).
	By the quadratic cancellation result upstream, the defect element
	\[
		\delta_{p_k(u)}^{(k)}(\square_{ij})
	\]
	lies in
	\[
		F^2K_{p_k(u)}^{(k)}.
	\]
	Its class modulo \(F^3\) is therefore well defined:
	\[
		\mathcal A_{ij}^{(k)}(p_k(u))
		:=
		\bigl[\delta_{p_k(u)}^{(k)}(\square_{ij})\bigr]
		\in
		F^2K_{p_k(u)}^{(k)}/F^3K_{p_k(u)}^{(k)}.
	\]
	Thus each state \(u\) determines, stage by stage, a family of
	finite-stage quadratic transport defect classes.
	\medskip
	\noindent
	\textbf{Step 2: Compatibility under refinement.}
	Let
	\[
		R_{k+1,k}^{\Omega}:
		\Omega_{p_{k+1}(u)}(\mathcal N_{k+1}(u))
		\to
		\Omega_{p_k(u)}(\mathcal N_k(u))
	\]
	and
	\[
		Q_{k+1,k}^{K}:
		K_{p_{k+1}(u)}^{(k+1)}
		\to
		K_{p_k(u)}^{(k)}
	\]
	be the bonding maps supplied by refinement compatibility.
	By the compatibility identity for defect maps,
	\[
		Q_{k+1,k}^{K}\circ \delta_{p_{k+1}(u)}^{(k+1)}
		=
		\delta_{p_k(u)}^{(k)}\circ R_{k+1,k}^{\Omega},
	\]
	the defect of the commutator square at stage \(k+1\) maps to the defect
	of the corresponding commutator square at stage \(k\).
	Passing to the quadratic quotient yields
	\[
		Q_{k+1,k}^{(2)}
		\bigl(\mathcal A_{ij}^{(k+1)}(p_{k+1}(u))\bigr)
		=
		\mathcal A_{ij}^{(k)}(p_k(u)),
	\]
	where
	\[
		Q_{k+1,k}^{(2)}:
		F^2K_{p_{k+1}(u)}^{(k+1)}/F^3K_{p_{k+1}(u)}^{(k+1)}
		\to
		F^2K_{p_k(u)}^{(k)}/F^3K_{p_k(u)}^{(k)}
	\]
	is the induced degree--\(2\) bonding morphism.
	Therefore
	\[
		\bigl(\mathcal A_{ij}^{(k)}(p_k(u))\bigr)_k
	\]
	is a compatible family in the inverse system defining the stabilized
	carrier.
	\medskip
	\noindent
	\textbf{Step 3: Passage to the inverse limit.}
	By definition,
	\[
		K=\mathcal K_\infty=\varprojlim_k \mathcal K_k.
	\]
	Hence every compatible family of finite-stage quadratic classes defines a
	unique element of \(K\).
	Applying this to the family attached to \(u\), we obtain an element
	\[
		\kappa(u)\in K.
	\]
	This constructs a map
	\[
		\kappa:U\to K,
		\qquad
		u\mapsto \kappa(u).
	\]
	At this stage the map is defined relative to the chosen basepoint system
	\((p_k(u))_k\). The next theorem removes that dependence.
	\medskip
	\noindent
	\textbf{Step 4: Finite-stage equivariance under intrinsic automorphisms.}
	Let
	\[
		g\in G=\Aut(U,\mathcal C).
	\]
	Because \(g\) preserves the comparison world, it preserves the induced
	transport structure. In particular, for each \(k\), it carries local
	witness systems at \(u\) to local witness systems at \(g\cdot u\), sends
	commutator squares to commutator squares, and induces an isomorphism
	\[
		g_*^{(k)}:
		F^2K_{p_k(u)}^{(k)}/F^3K_{p_k(u)}^{(k)}
		\to
		F^2K_{p_k(g\cdot u)}^{(k)}/F^3K_{p_k(g\cdot u)}^{(k)}
	\]
	satisfying
	\[
		g_*^{(k)}
		\bigl(\mathcal A_{ij}^{(k)}(p_k(u))\bigr)
		=
		\mathcal A_{ij}^{(k)}(p_k(g\cdot u)).
	\]
	Since the family \((g_*^{(k)})_k\) is compatible with refinement, it
	assembles into an induced automorphism
	\[
		g_*:K\to K.
	\]
	\medskip
	\noindent
	\textbf{Step 5: Naturality.}
	The \(k\)-th component of \(\kappa(u)\) is
	\[
		\mathcal A_{ij}^{(k)}(p_k(u)).
	\]
	Applying \(g_*\) to \(\kappa(u)\) acts componentwise by \(g_*^{(k)}\), so
	the \(k\)-th component of \(g_*\kappa(u)\) is
	\[
		g_*^{(k)}
		\bigl(\mathcal A_{ij}^{(k)}(p_k(u))\bigr)
		=
		\mathcal A_{ij}^{(k)}(p_k(g\cdot u)).
	\]
	But this is exactly the \(k\)-th component of \(\kappa(g\cdot u)\).
	Since equality in an inverse limit is equivalent to equality of all
	finite-stage components, it follows that
	\[
		\kappa(g\cdot u)=g_*\kappa(u).
	\]
	\medskip
	\noindent
	\textbf{Step 6: Conclusion.}
	Thus the map
	\[
		\kappa:U\to K
	\]
	is canonically defined from compatible quadratic defect classes and is
	natural under the intrinsic symmetry action.
\end{proof}
\begin{theorem}[Basepoint-independence of the quadratic defect section]
	Let the hypotheses of \cref{thm:fusion-kappa} hold.
	Fix a state
	\[
		u\in U.
	\]
	At each finite refinement stage \(k\), let
	\[
		p_k,\;p_k'\in V_k(u)
	\]
	be two admissible basepoints representing the same state \(u\), and let
	\[
		\gamma_k:p_k\to p_k'
	\]
	be an admissible transport path between them.
	Let
	\[
		\mathcal A_{ij}^{(k)}(p_k),
		\qquad
		\mathcal A_{ij}^{(k)}(p_k')
	\]
	be the quadratic defect classes computed using these two basepoints.
	Then
	\[
		\mathcal A_{ij}^{(k)}(p_k')
		=
		\mathcal A_{ij}^{(k)}(p_k)
		\quad\text{in}\quad
		F^2K^{(k)}/F^3K^{(k)}.
	\]
	Consequently, the element
	\[
		\kappa(u)\in K
	\]
	is independent of the choice of basepoint representatives.
\end{theorem}
\begin{proof}
	We prove the theorem in four steps.
	\medskip
	\noindent
	\textbf{Step 1: Conjugation identifies based loop groups.}
	Fix a stage \(k\).
	Let
	\[
		\gamma_k:p_k\to p_k'
	\]
	be an admissible transport path.
	Then conjugation by \(\gamma_k\) identifies the based loop groups:
	\[
		\Omega_{p_k}(\mathcal N_k(u))
		\longrightarrow
		\Omega_{p_k'}(\mathcal N_k(u)),
		\qquad
		\ell\longmapsto \gamma_k\,\ell\,\gamma_k^{-1}.
	\]
	Under the transport functor, this induces conjugation on the isotropy
	group:
	\[
		\delta_{p_k'}^{(k)}(\gamma_k\ell\gamma_k^{-1})
		=
		\delta_{\gamma_k}^{(k)}\,
		\delta_{p_k}^{(k)}(\ell)\,
		\delta_{\gamma_k}^{(k)\,-1},
	\]
	where
	\[
		\delta_{\gamma_k}^{(k)}\in K^{(k)}
	\]
	denotes the transport element associated to \(\gamma_k\).
	\medskip
	\noindent
	\textbf{Step 2: The two commutator-square defects are conjugate.}
	Let
	\[
		\square_{ij}:=e_ie_je_i^{-1}e_j^{-1}
	\]
	be the commutator square based at \(p_k\), and let
	\[
		\square_{ij}':=\gamma_k\,\square_{ij}\,\gamma_k^{-1}
	\]
	be the corresponding square at \(p_k'\).
	Then
	\[
		\delta_{p_k'}^{(k)}(\square_{ij}')
		=
		\delta_{\gamma_k}^{(k)}\,
		\delta_{p_k}^{(k)}(\square_{ij})\,
		\delta_{\gamma_k}^{(k)\,-1}.
	\]
	Set
	\[
		x:=\delta_{p_k}^{(k)}(\square_{ij})\in F^2K^{(k)},
		\qquad
		g:=\delta_{\gamma_k}^{(k)}\in F^1K^{(k)}.
	\]
	Then the two representatives are related by
	\[
		\delta_{p_k'}^{(k)}(\square_{ij}')=gxg^{-1}.
	\]
	\medskip
	\noindent
	\textbf{Step 3: Conjugation is trivial modulo \(F^3\) on \(F^2\).}
	Using the identity
	\[
		gxg^{-1}=x[g,x],
	\]
	we reduce the difference between \(gxg^{-1}\) and \(x\) to the
	commutator term \([g,x]\).
	Since
	\[
		g\in F^1,\qquad x\in F^2,
	\]
	\cref{lem:fusion-commutator} yields
	\[
		[g,x]\in [F^1,F^2]\subseteq F^3.
	\]
	Hence
	\[
		gxg^{-1}=x\cdot z
		\qquad
		\text{for some }z\in F^3.
	\]
	Therefore
	\[
		gxg^{-1}\equiv x\pmod{F^3}.
	\]
	It follows that
	\[
		\bigl[\delta_{p_k'}^{(k)}(\square_{ij}')\bigr]
		=
		\bigl[\delta_{p_k}^{(k)}(\square_{ij})\bigr]
		\quad\text{in}\quad
		F^2K^{(k)}/F^3K^{(k)}.
	\]
	Equivalently,
	\[
		\mathcal A_{ij}^{(k)}(p_k')
		=
		\mathcal A_{ij}^{(k)}(p_k).
	\]
	\medskip
	\noindent
	\textbf{Step 4: Passage to the inverse limit.}
	Since the equality above holds for every finite stage \(k\), the two
	compatible families obtained from the two basepoint systems coincide in
	every component. Hence they define the same element of the inverse limit
	\[
		K=\varprojlim_k \mathcal K_k.
	\]
	Therefore the stabilized defect class
	\[
		\kappa(u)\in K
	\]
	is independent of the choice of admissible basepoint representatives.
\end{proof}
\begin{theorem}[Functorial uniqueness of the quadratic defect section]
	Let the hypotheses of \cref{thm:fusion-kappa} hold, and let
	\[
		\kappa:U\to K
	\]
	be the canonical quadratic defect section just constructed.
	Suppose
	\[
		\sigma:U\to K
	\]
	is another map satisfying the following properties:
	\begin{enumerate}[label=\textup{(\roman*)}]
		\item
		      for every
		      \[
			      g\in G=\Aut(U,\mathcal C)
			      \qquad\text{and}\qquad
			      u\in U,
		      \]
		      one has
		      \[
			      \sigma(g\cdot u)=g_*\sigma(u);
		      \]
		\item
		      for each
		      \[
			      u\in U,
		      \]
		      the value \(\sigma(u)\) depends only on the first visible quadratic
		      transport defect class of \(u\), i.e. only on the corresponding element
		      of the stabilized carrier \(K\);
		\item
		      on primitive quadratic square witnesses, \(\sigma\) agrees with the
		      finite-stage defect assignment used to define \(\kappa\).
	\end{enumerate}
	Then
	\[
		\sigma=\kappa.
	\]
	That is, \(\kappa\) is the unique natural quadratic defect assignment
	compatible with transport, filtration, and primitive quadratic
	normalization.
\end{theorem}
\begin{proof}
	We prove equality pointwise.
	Fix
	\[
		u\in U.
	\]
	By \cref{thm:fusion-kappa}, the element
	\[
		\kappa(u)\in K
	\]
	is represented by a compatible family of finite-stage quadratic square
	classes
	\[
		\bigl(\mathcal A_{ij}^{(k)}(p_k(u))\bigr)_k.
	\]
	By construction of the stabilized carrier, \(K\) is generated by the
	stabilized images of such primitive quadratic square classes.
	Now consider
	\[
		\sigma(u)\in K.
	\]
	Property \textup{(ii)} says that \(\sigma(u)\) depends only on the first
	visible quadratic transport defect class of \(u\). But this class is
	exactly the stabilized class represented by the same compatible family
	used to define \(\kappa(u)\). Thus \(\sigma(u)\) cannot depend on any
	additional structure beyond the element of \(K\) already encoded by
	\(\kappa(u)\).
	Property \textup{(iii)} fixes the normalization of \(\sigma\) on the
	primitive quadratic generators: on every primitive quadratic square
	witness, \(\sigma\) agrees with the defining assignment for \(\kappa\).
	Since the stabilized carrier \(K\) is generated by these primitive
	quadratic classes, and since \(\sigma\) and \(\kappa\) agree on those
	generators while depending only on the resulting stabilized quadratic
	class, it follows that
	\[
		\sigma(u)=\kappa(u).
	\]
	Because \(u\in U\) was arbitrary, we conclude that
	\[
		\sigma=\kappa.
	\]
	This proves uniqueness.
\end{proof}
% ============================================================
\section{Construction of the invariant}
\label{sec:fusion-construction}
% ============================================================
We now construct the confinement invariant itself.
\begin{definition}[Quadratic defect section]
	Let
	\[
		\kappa:U\to K
	\]
	be the canonical degree--\(2\) defect section of
	\cref{thm:fusion-kappa}.
\end{definition}
\begin{definition}[The confinement invariant]
	Define
	\[
		\mathfrak s_{\mathrm{fus}}:=Q\circ \kappa.
	\]
\end{definition}
\begin{theorem}[Existence]
	\label{thm:fusion-existence}
	The map
	\[
		\mathfrak s_{\mathrm{fus}}:U\to\mathbb R
	\]
	is a well-defined scalar-valued function on the confinement state space.
\end{theorem}
\begin{proof}
	To prove that \(\mathfrak s_{\mathrm{fus}}\) is well defined, it suffices to verify that both
	factors in the composite
	\[
		\mathfrak s_{\mathrm{fus}}=Q\circ\kappa
	\]
	are well-defined maps with compatible source and target.
	First, by \cref{thm:fusion-kappa}, the quadratic defect section
	\[
		\kappa:U\to K
	\]
	is a canonically defined map on the confinement comparison world. The
	word ``canonical'' is essential: it means that \(\kappa(u)\) depends
	only on the intrinsic transport data associated with the state \(u\), and
	not on any auxiliary presentation, lift, or coordinate choice. Thus the
	expression \(\kappa(u)\) makes sense for every \(u\in U\).
	Second, by \cref{prop:fusion-scalar}, the scalar channel
	\[
		Q:K\to\mathbb R
	\]
	is a well-defined real-valued map on the stabilized quadratic carrier.
	Again, this means that once a degree--\(2\) transport defect class has
	been formed as an element of \(K\), its scalar image under \(Q\) is
	unambiguously determined.
	Therefore, for each \(u\in U\), one may first form
	\[
		\kappa(u)\in K
	\]
	and then apply \(Q\) to obtain
	\[
		Q(\kappa(u))\in\mathbb R.
	\]
	This defines a map
	\[
		\mathfrak s_{\mathrm{fus}}:U\to\mathbb R,
		\qquad
		\mathfrak s_{\mathrm{fus}}(u):=Q(\kappa(u)).
	\]
	No further compatibility condition is required, because the codomain of
	\(\kappa\) is precisely the domain of \(Q\).
	Hence \(\mathfrak s_{\mathrm{fus}}\) is a well-defined scalar-valued function on \(U\).
\end{proof}
\begin{theorem}[Invariance under intrinsic evolution]
	\label{thm:fusion-invariance}
	Under the standing principle \cref{sp:closed-world}, let
	\[
		\Phi_t\in \Aut(U,\mathcal C)
	\]
	be an intrinsic evolution of the confinement system.
	Then
	\[
		\mathfrak s_{\mathrm{fus}}(\Phi_t(u))=\mathfrak s_{\mathrm{fus}}(u)
		\qquad
		\text{for all }u\in U,\ t.
	\]
\end{theorem}
\begin{proof}
	We prove the statement by unfolding, in order, the definitions of
	intrinsic evolution, naturality of the defect section, and admissibility
	of the scalar channel.
	Fix \(t\) and \(u\in U\).
	By hypothesis,
	\[
		\Phi_t\in \Aut(U,\mathcal C).
	\]
	Thus \(\Phi_t\) is an automorphism of the closed comparison world.
	Equivalently, \(\Phi_t\) preserves all admissible comparison predicates
	and therefore belongs to the intrinsic symmetry group
	\[
		G:=\Aut(U,\mathcal C).
	\]
	This is consistent with \cref{thm:transport-forces-Aut}, which asserts
	precisely that intrinsic admissible transport-closed dynamics acts
	through automorphisms of the comparison world.
	Now apply the functoriality of the defect section.
	By \cref{thm:fusion-kappa}, for every \(g\in G\),
	\[
		\kappa(g\cdot u)=g_*\kappa(u),
	\]
	where \(g_*\) denotes the induced action on the stabilized quadratic
	carrier \(K\).
	Specializing to \(g=\Phi_t\), we obtain
	\[
		\kappa(\Phi_t(u))=\Phi_{t*}\kappa(u).
	\]
	Next use the admissibility of the scalar channel \(Q\).
	By definition of admissible scalarization, and by the uniqueness theory
	cited in \cref{prop:fusion-scalar}, the map \(Q\) is invariant under the
	intrinsic symmetry action on \(K\). That is,
	\[
		Q(g_*k)=Q(k)
		\qquad
		\text{for all }g\in G,\ k\in K.
	\]
	Again specializing to \(g=\Phi_t\), we obtain
	\[
		Q(\Phi_{t*}k)=Q(k)
		\qquad
		\text{for all }k\in K.
	\]
	We now compute directly:
	\begin{align*}
		\mathfrak s_{\mathrm{fus}}(\Phi_t(u))
		 & =(Q\circ\kappa)(\Phi_t(u)) \\
		 & =Q(\kappa(\Phi_t(u)))      \\
		 & =Q(\Phi_{t*}\kappa(u))     \\
		 & =Q(\kappa(u))              \\
		 & =\mathfrak s_{\mathrm{fus}}(u).
	\end{align*}
	Since \(u\in U\) and \(t\) were arbitrary, the equality holds for all
	states and all intrinsic evolution times.
	Therefore \(\mathfrak s_{\mathrm{fus}}\) is invariant under intrinsic evolution.
\end{proof}
\begin{theorem}[Uniqueness up to scale]
	\label{thm:fusion-unique}
	The scalar invariant \(\mathfrak s_{\mathrm{fus}}\) is unique up to multiplication by a nonzero
	real scalar.
\end{theorem}
\begin{proof}
	Let
	\[
		\mathfrak s_1,\mathfrak s_2:U\to\mathbb R
	\]
	be two scalar invariants obtained by the construction of the appendix.
	By definition of the construction, there exist nonzero admissible
	scalarizations
	\[
		Q_1,Q_2:K\to\mathbb R
	\]
	such that
	\[
		\mathfrak s_1=Q_1\circ\kappa,
		\qquad
		\mathfrak s_2=Q_2\circ\kappa.
	\]
	By \cref{prop:fusion-scalar}, the admissible scalarization space on \(K\)
	is one-dimensional.
	Therefore there exists
	\[
		c\in\mathbb R^\times
	\]
	such that
	\[
		Q_1=c\,Q_2.
	\]
	Substituting this relation into the definition of \(\mathfrak s_1\), we obtain
	\[
		\mathfrak s_1
		=
		Q_1\circ\kappa
		=
		(c\,Q_2)\circ\kappa
		=
		c\,(Q_2\circ\kappa)
		=
		c\,\mathfrak s_2.
	\]
	Thus \(\mathfrak s_1\) and \(\mathfrak s_2\) differ by multiplication by a nonzero real
	constant.
	This proves uniqueness of the confinement invariant up to nonzero
	normalization.
\end{proof}
\begin{corollary}[Main appendix theorem]
	\label{cor:fusion-main}
	The confinement comparison world carries a scalar-valued function
	\[
		\mathfrak s_{\mathrm{fus}}:U\to\mathbb R
	\]
	which is invariant under intrinsic evolution and unique up to
	multiplication by a nonzero scalar.
\end{corollary}
\begin{proof}
	The corollary is obtained by assembling
	\cref{thm:fusion-existence,thm:fusion-invariance,thm:fusion-unique}.
	By \cref{thm:fusion-existence}, the function
	\[
		\mathfrak s_{\mathrm{fus}}=Q\circ\kappa
	\]
	is well defined on the confinement state space \(U\).
	By \cref{thm:fusion-invariance}, this function is preserved under every
	intrinsic evolution
	\[
		\Phi_t\in\Aut(U,\mathcal C).
	\]
	By \cref{thm:fusion-unique}, any other scalar invariant obtained from the
	same transport stack differs from \(\mathfrak s_{\mathrm{fus}}\) by multiplication by a nonzero
	real scalar.
	These three statements are exactly the claims asserted in the corollary.
\end{proof}
% ============================================================
\section{Perturbative alignment criterion}
\label{sec:fusion-asymptotic}
% ============================================================
\Cref{cor:fusion-main} is structural and nonperturbative.
We now record the additional criterion under which one may identify it
with the perturbative adiabatic invariant in the small-\(\epsilon\)
regime.
\begin{theorem}[Conditional carrier identification of the perturbative invariant]
	\label{thm:fusion-equivalence}
	Under the standing principle \cref{sp:closed-world}, assume that in the
	perturbative regime
	\[
		\epsilon:=\rho/L\ll 1,
	\]
	the perturbative adiabatic invariant is obtained by admissible
	scalarization of the same transport stack, with leading nontrivial term
	carried by the first visible transport defect.
	Then the perturbative adiabatic invariant and the closed-system invariant
	\[
		\mathfrak s_{\mathrm{fus}}=Q\circ\kappa
	\]
	are the same scalar channel written in different coordinates.
	Equivalently, under that perturbative identification, Kruskal's
	adiabatic invariant series is the asymptotic expansion of the unique
	admissible scalar channel on the first visible transport carrier.
\end{theorem}
\begin{proof}
	The proof rests on four assertions:
	\begin{enumerate}[label=\textup{(\arabic*)}]
		\item the transport filtration records perturbative order;
		\item the first visible scalar carrier occurs in degree \(2\);
		\item the scalar channel on that carrier is unique;
		\item under the stated perturbative hypothesis, no second independent
		      scalar channel enters the same transport stack.
	\end{enumerate}
	We treat them in order and then combine them.
	\medskip
	\noindent
	\textbf{Step 1: The filtration records perturbative order.}
	The small parameter
	\[
		\epsilon:=\rho/L
	\]
	measures the ratio of the microscopic gyroradius scale to the macroscopic
	variation scale of the magnetic geometry.
	In the perturbative treatment of magnetic confinement, transport
	quantities are expanded as formal series in \(\epsilon\), and each order
	measures the first scale at which exact closure of the leading transport
	law fails \cite{Northrop1963}.
	The intrinsic transport filtration
	\[
		F^1\supset F^2\supset F^3\supset\cdots
	\]
	is organized by exactly the same principle.
	An element lies in \(F^m\) precisely when its transport defect is
	invisible through order \(m-1\), and passage to the quotient
	\[
		F^m/F^{m+1}
	\]
	isolates the first nonvanishing defect at order \(m\).
	Thus the perturbative order-\(\epsilon^m\) data are encoded by the graded
	piece
	\[
		F^m/F^{m+1}.
	\]
	\medskip
	\noindent
	\textbf{Step 2: The first visible scalar carrier is quadratic.}
	By the main transport theory and, in particular, by the interface
	results
	(\cref{thm:interface-stabilized,thm:interface-seal}), the first visible nontrivial transport
	obstruction is degree \(2\).
	Equivalently, the first graded piece capable of carrying nontrivial
	scalar information is
	\[
		K\cong F^2/F^3.
	\]
	The degree--\(1\) quotient carries only abelianized transport content
	and does not detect the intrinsic \(2\)-skeletal obstruction.
	Hence the first nontrivial scalar invariant extracted from transport must
	factor through \(K\).
	\medskip
	\noindent
	\textbf{Step 3: Uniqueness of the scalar channel on the first visible carrier.}
	By \cref{prop:fusion-scalar}, the space of admissible scalarizations of
	\(K\) in the realized scalar-channel regime is one-dimensional.
	Therefore any admissible scalar observable whose leading contribution is
	carried by the first visible transport defect must be proportional to
	the distinguished scalar channel
	\[
		Q:K\to\mathbb R.
	\]
	By the perturbative-identification hypothesis, the perturbative
	adiabatic invariant is of exactly this type.
	\medskip
	\noindent
	\textbf{Step 4: Identification.}
	Under the same perturbative-identification hypothesis, higher-order
	corrections still scalarize the same transport stack rather than opening
	a new independent scalar channel.
	By \cref{thm:kernel-annihilation-upstream}, every admissible degree--\(2\)
	scalar detector descends uniquely through the quadratic carrier.
	Thus once the leading scalar transport content has been fixed on \(K\),
	there is no second independent scalar channel through which that same
	transport stack can be scalarized.
	Consequently, the perturbative invariant and the intrinsic invariant
	\[
		\mathfrak s_{\mathrm{fus}}=Q\circ\kappa
	\]
	arise from the same scalar channel and differ only by the perturbative
	coordinate expansion used to express that channel in the regime
	\(\epsilon\ll1\).
	This is the asserted identification.
\end{proof}
\begin{corollary}[Conditional all-orders asymptotic agreement]
	Under the hypotheses of \cref{thm:fusion-equivalence}, in the
	perturbative regime
	\[
		\epsilon\ll 1,
	\]
	the invariant \(\mathfrak s_{\mathrm{fus}}\) agrees with Kruskal's adiabatic invariant series to
	all orders.
\end{corollary}
\begin{proof}
	By \cref{thm:fusion-equivalence}, under the stated perturbative
	identification the perturbative invariant series and the intrinsic
	invariant \(\mathfrak s_{\mathrm{fus}}\) are two presentations of the same scalar channel.
	Therefore their asymptotic coefficients agree order by order in the
	small-\(\epsilon\) regime.
\end{proof}
\begin{remark}
	The point is not that \(\mathfrak s_{\mathrm{fus}}\) is defined perturbatively.
	It is not.
	Rather, once \(\mathfrak s_{\mathrm{fus}}\) has been defined intrinsically by the closed-system
	transport stack, any perturbative invariant satisfying
	\cref{thm:fusion-equivalence} is recovered as a local coordinate
	expansion of that same scalar channel in the regime
	\(\epsilon\ll1\).
\end{remark}
% ============================================================
\section{Conditional confinement package}
\label{sec:fusion-prediction}
% ============================================================
The theorem-level output of the appendix is the intrinsic scalar \(\mathfrak s_{\mathrm{fus}}\).
We now record the perturbative and confinement readings that become
available once the relevant interpretation hypotheses are fixed.
\begin{theorem}[Conditional confinement interpretation package]
	\label{thm:fusion-prediction}
	Under the standing principle \cref{sp:closed-world}, let
	\((U,\mathcal C)\) be a magnetic geometry realized as a closed
	comparison world.
	Then the following statements hold.
	\begin{enumerate}[label=\textup{(\roman*)}]
		\item
		      There exists a scalar-valued function
		      \[
			      \mathfrak s_{\mathrm{fus}}:U\to\mathbb R
		      \]
		      which is invariant under intrinsic evolution and unique up to
		      multiplication by a nonzero real scalar.
		\item
		      Under the additional perturbative-identification hypothesis of
		      \cref{thm:fusion-equivalence}, the perturbative adiabatic invariant is the
		      asymptotic expansion of \(\mathfrak s_{\mathrm{fus}}\) in the regime \(\epsilon\ll 1\).
		\item
		      If confinement for the relevant particle family is modeled by compact
		      containment of the appropriate invariant level sets, then the resulting
		      structural confinement criterion is that the magnetic geometry confines
		      exactly when those level sets of \(\mathfrak s_{\mathrm{fus}}\) are compact and remain contained in
		      the plasma volume.
		\item
		      Under the same confinement model, the corresponding optimization problem is
		      to optimize the confinement quality of the level sets of \(\mathfrak s_{\mathrm{fus}}\) over the
		      admissible comparison structures induced by the geometry.
	\end{enumerate}
\end{theorem}
\begin{proof}
	Part \textup{(i)} is exactly \cref{cor:fusion-main}.
	Part \textup{(ii)} is exactly \cref{thm:fusion-equivalence}.
	For parts \textup{(iii)} and \textup{(iv)}, once confinement is modeled by compact
	containment of invariant level sets, the statements are the formal
	restatement of that model using the theorem-level scalar \(\mathfrak s_{\mathrm{fus}}\) singled
	out by part \textup{(i)}.
	Because changing magnetic geometry changes the induced comparison world,
	and hence changes the resulting invariant \(\mathfrak s_{\mathrm{fus}}\), optimizing confinement
	inside that model is equivalently the problem of optimizing the geometry
	of the relevant \(\mathfrak s_{\mathrm{fus}}\)-level sets across the induced comparison
	structures.
\end{proof}
\begin{remark}
	The theorem does not yet provide a reactor-specific closed-form number
	for \(\mathfrak s_{\mathrm{fus}}\) for an arbitrary coil set.
	What it provides is the canonical intrinsic scalar to which any later
	perturbative, numerical, or confinement model must refer.
\end{remark}
% ============================================================
\section{Magnetic-confinement interpretation}
\label{sec:fusion-interpretation}
% ============================================================
\Cref{thm:fusion-prediction} separates theorem-level invariant content
from the later confinement reading.
Under the confinement interpretation package, the structure may be read
as follows; the confinement terminology is the standard magnetic
confinement context rather than an additional theorem of the appendix
\cite{Wilson1974}.
\begin{enumerate}[label=\textup{(\roman*)}]
	\item
	      The gyration of a charged particle in a magnetic field is an internal
	      transport process of the confinement world.
	\item
	      The failure of this transport to close exactly is a loop defect.
	      Its first stable visible component lies in the quadratic carrier
	      \[
		      F^2/F^3.
	      \]
	\item
	      The map
	      \[
		      \kappa:U\to K
	      \]
	      assigns to each state its intrinsic quadratic transport defect class.
	\item
	      The scalar channel
	      \[
		      Q:K\to\mathbb R
	      \]
	      extracts the unique admissible scalar carried by that defect in the
	      realized scalar-channel regime used in the appendix.
	\item
	      The composite
	      \[
		      \mathfrak s_{\mathrm{fus}}=Q\circ\kappa
	      \]
	      is therefore the canonical intrinsic transport scalar supplied by
	      the closed-system structure for this application.
\end{enumerate}
Thus the intrinsic transport scalar is not postulated at the
carrier level.
Its later reading as a nonperturbative adiabatic invariant or
confinement observable is a downstream interpretation of the scalar
channel of the quadratic transport carrier.
\begin{remark}
	No Liouville-integrability assumption enters.
	No continuous-symmetry assumption enters.
	No data-driven learning step enters.
	All three are replaced at the carrier level by the closed-system
	transport stack already proved in the main text.
\end{remark}
% ============================================================
\section{Consequences of the interpretation package}
\label{sec:fusion-predictions}
% ============================================================
Once the intrinsic scalar \(\mathfrak s_{\mathrm{fus}}\) is established theorem-level and the
perturbative/level-set interpretation package is adopted, several
immediate consequences follow.
\subsection{General-geometry existence}
The existence of \(\mathfrak s_{\mathrm{fus}}\) is not restricted to integrable or symmetric
magnetic fields.
Its construction depends only on the closed comparison structure and the
quadratic scalar channel.
Accordingly, any magnetic geometry that can be realized as a closed
comparison world carries the same canonical nonperturbative transport
scalar.
\subsection{Structural confinement criterion under level-set interpretation}
Because \(\mathfrak s_{\mathrm{fus}}\) is invariant under intrinsic evolution, if confinement is
modeled by compact containment of the relevant invariant level sets, then
confinement may be rephrased geometrically in terms of the level sets of
\(\mathfrak s_{\mathrm{fus}}\).
The resulting structural criterion is:
\begin{quote}
	A magnetic geometry confines the relevant particles exactly when the
	appropriate level sets of \(\mathfrak s_{\mathrm{fus}}\) are compact and remain contained within
	the plasma volume.
\end{quote}
This criterion is downstream: it depends on the theorem-level scalar
channel of the transport filtration together with the chosen level-set
interpretation, not on full-orbit simulation.
\subsection{Optimization reformulation under the same criterion}
Changing magnetic geometry changes the comparison predicates.
Changing the comparison predicates changes the transport filtration.
Changing the transport filtration changes the quadratic carrier and
therefore changes the scalar channel \(\mathfrak s_{\mathrm{fus}}\).
Thus, under the same level-set confinement model, the optimization
problem becomes:
\begin{quote}
	Find the magnetic geometry whose induced comparison structure optimizes
	the compactness and confinement properties of the level sets of \(\mathfrak s_{\mathrm{fus}}\).
\end{quote}
This reformulates the task as a variational problem on the space of
comparison structures rather than only as a brute-force orbit-search
problem.
\subsection{Further confinement regimes}
Nothing in the theorem-level construction of \(\mathfrak s_{\mathrm{fus}}\) used a special
property of alpha particles beyond the existence of a closed comparison
structure for the confinement system.
Accordingly, the same structural machinery applies, in the same
conditional sense, to any confinement problem whose relevant state space
is describable as a closed comparison world with transport filtration
and quadratic scalar channel.
% ============================================================
\section{Normalized exact observables}
\label{sec:fusion-normalized}
% ============================================================
The scalar channel \(Q\) is unique only up to multiplication by a
nonzero constant.
Accordingly, absolute values of \(\mathfrak s_{\mathrm{fus}}=Q\circ\kappa\) depend on the chosen
normalization.
What is canonically determined without any further choice is the ratio of
two nonzero scalar values.
\begin{proposition}[Normalized exact invariant]
	Let
	\[
		Q:K\to\mathbb R_{\ge 0}
	\]
	be the canonical scalar channel on the stabilized quadratic carrier, and
	let
	\[
		u_*\in U
	\]
	satisfy
	\[
		\kappa(u_*)\neq 0.
	\]
	Then the normalized scalar
	\[
		\widehat{\mathfrak s}_{\mathrm{fus}}(u;u_*):=
		\frac{Q(\kappa(u))}{Q(\kappa(u_*))}
	\]
	is well defined, independent of the normalization of \(Q\), and
	invariant under intrinsic evolution.
\end{proposition}
\begin{proof}
	There are three claims to establish:
	\begin{enumerate}[label=\textup{(\roman*)}]
		\item the ratio is independent of normalization of \(Q\);
		\item the ratio is defined, i.e. the denominator is nonzero;
		\item the ratio is invariant under intrinsic evolution.
	\end{enumerate}
	\medskip
	\noindent
	\textbf{Step 1: Independence of normalization.}
	Let
	\[
		Q'=cQ
	\]
	for some \(c>0\).
	Then for every \(u\in U\),
	\[
		Q'(\kappa(u))=c\,Q(\kappa(u)),
	\]
	and similarly
	\[
		Q'(\kappa(u_*))=c\,Q(\kappa(u_*)).
	\]
	Therefore
	\[
		\frac{Q'(\kappa(u))}{Q'(\kappa(u_*))}
		=
		\frac{cQ(\kappa(u))}{cQ(\kappa(u_*))}
		=
		\frac{Q(\kappa(u))}{Q(\kappa(u_*))}.
	\]
	Thus the normalized scalar does not depend on the chosen nonzero
	normalization of the scalar channel.
	\medskip
	\noindent
	\textbf{Step 2: Nonvanishing of the denominator.}
	By hypothesis,
	\[
		\kappa(u_*)\neq 0.
	\]
	The scalar channel \(Q\) is by construction nontrivial on the realized
	carrier and takes nonnegative values on the chosen normalization.
	The intended normalization in this proposition is one for which nonzero
	carrier classes are assigned nonzero nonnegative values. Under that
	normalization,
	\[
		Q(\kappa(u_*))>0.
	\]
	Hence the denominator is nonzero and the ratio is well defined.
	\medskip
	\noindent
	\textbf{Step 3: Invariance under intrinsic evolution.}
	Let
	\[
		\Phi_t\in \Aut(U,\mathcal C).
	\]
	By \cref{thm:fusion-invariance},
	\[
		Q(\kappa(\Phi_t(u)))=Q(\kappa(u))
	\]
	for every \(u\in U\).
	Applying this once to \(u\) and once to \(u_*\), we obtain
	\[
		Q(\kappa(\Phi_t(u)))=Q(\kappa(u)),
		\qquad
		Q(\kappa(\Phi_t(u_*)))=Q(\kappa(u_*)).
	\]
	Therefore
	\[
		\widehat{\mathfrak s}_{\mathrm{fus}}(\Phi_t(u);\Phi_t(u_*))
		=
		\frac{Q(\kappa(\Phi_t(u)))}{Q(\kappa(\Phi_t(u_*)))}
		=
		\frac{Q(\kappa(u))}{Q(\kappa(u_*))}
		=
		\widehat{\mathfrak s}_{\mathrm{fus}}(u;u_*).
	\]
	This proves invariance of the normalized scalar.
\end{proof}
\begin{remark}
	The proposition identifies the first exactly computable scalar quantity
	determined entirely by the theorem stack without fixing an external
	normalization convention.
\end{remark}
% ============================================================
\section{Structural summary}
\label{sec:fusion-summary}
% ============================================================
The argument of the appendix is:
\begin{align*}
	\text{closed confinement system} & \Longrightarrow \text{comparison completeness}                                      \\
	                                 & \Longrightarrow \text{canonical quotient semantics}                                 \\
	                                 & \Longrightarrow \text{transport filtration } F^1\supset F^2\supset F^3\supset\cdots \\
	                                 & \Longrightarrow \text{stabilized quadratic carrier } K                              \\
	                                 & \Longrightarrow \text{unique scalar channel } Q                                     \\
	                                 & \Longrightarrow \text{intrinsic scalar } \mathfrak s_{\mathrm{fus}} .
\end{align*}
Through the intrinsic scalar \(\mathfrak s_{\mathrm{fus}}=Q\circ\kappa\), every arrow in this
chain is theorem-level within the closed relational stack, using the
inherited realized scalar-channel package already isolated upstream.
Thus \(\mathfrak s_{\mathrm{fus}}\) is not postulated, not learned from data, and not restricted
to integrable geometries at the level of its intrinsic construction.
Perturbative matching and confinement-level interpretation are
additional downstream reading layers built on that scalar.
% ============================================================
\section{Conclusion}
\label{sec:fusion-conclusion}
% ============================================================
This appendix has one theorem-level purpose:
to exhibit the downstream confinement use of the quadratic scalar channel
of the closed-system transport theory.
The chain proved here is:
\begin{align*}
	(U,\mathcal{C}) & \Longrightarrow N=F^1                                \\
	                & \Longrightarrow V=F^1/F^2                            \\
	                & \Longrightarrow \Lambda^2(V) \xrightarrow{q_{\wedge}} F^2/F^3 \\
	                & \Longrightarrow K = \mathcal{K}_\infty               \\
	                & \Longrightarrow Q : K \to \mathbb{R}                 \\
	                & \Longrightarrow \mathfrak s_{\mathrm{fus}} = Q \circ \kappa .
\end{align*}
Thus the abstract machinery of the main text yields a canonical intrinsic
transport scalar in the magnetic-confinement setting whenever the system is
realized as a closed comparison world.  Its existence and uniqueness up to
scale are fixed by the internal structure of that world; perturbative
matching to Kruskal-type series and any confinement-level or optimization
reading remain downstream interpretation steps.

\begin{chapterclosing}
\closingitem{Established}{The closed transport stack yields a canonical intrinsic scalar for confinement systems realized as closed comparison worlds.}
\closingitem{Carried forward}{Perturbative matching and confinement-level interpretation remain downstream readings of that scalar.}
\end{chapterclosing}
